\section*{Overview}

Monotone queue optimization is a DP speedup for transitions that ask for the minimum or maximum over a sliding window
of previous states. Instead of scanning that window every time, a deque keeps only the candidates that can still win in
the future.

\section*{When to Use It}

Use it when:

\begin{itemize}
  \item the transition is \(\min\) or \(\max\) over a contiguous range of previous indices,
  \item the valid range slides forward as the current index increases,
  \item the value compared inside the transition can be maintained in monotone order.
\end{itemize}

\section*{Core Idea}

Suppose
\[
dp[i] = a[i] + \min_{i-k \le j < i} dp[j].
\]
As \(i\) moves right, the allowed range of \(j\) also moves right. A deque can store candidate indices in increasing
order of their DP values while discarding:
\begin{itemize}
  \item indices that leave the range,
  \item indices whose value is dominated by a newer candidate.
\end{itemize}

\section*{Key Insight}

If a newer index has a DP value no worse than an older one, and it will remain valid at least as long, the older one
can never become optimal again. That is what makes the deque small.

\section*{Operations / Main Technique}

\begin{itemize}
  \item pop from the front while an index is out of range,
  \item use the front as the best candidate,
  \item pop from the back while the new value dominates older candidates,
  \item push the current index.
\end{itemize}

\paragraph{Worked example.}
If the deque stores candidate DP values \([5, 7, 9]\) and the new state has value \(6\), then \(7\) and \(9\) are
removed from the back because they will never beat \(6\) in any future window where \(6\) is still present.

\section*{Correctness Intuition}

The deque remains ordered by increasing candidate value, so its front is always the best valid transition source. The
domination rule is safe because any removed candidate is worse than a newer one and expires no later.

\section*{Complexity Analysis}

Each index enters and leaves the deque at most once, so the optimized transition runs in \(O(n)\) total time.

\section*{Implementation}

The sample code solves the sliding-window minimum DP pattern directly. In practice, many problems reduce to that form
after algebraic cleanup of the transition.

\section*{Common Pitfalls}

\begin{itemize}
  \item Applying the trick when the valid transition range is not contiguous.
  \item Forgetting which transformed value should be monotone in the deque.
  \item Off-by-one mistakes in the window endpoints.
\end{itemize}

\section*{Variants / Extensions}

\begin{itemize}
  \item Sliding-window maximum by reversing the comparisons.
  \item Deque optimization after subtracting a linear term from the transition.
  \item Comparison with convex hull trick when the dependence is linear instead of window-based.
\end{itemize}

\section*{Practice Problems}

\begin{itemize}
  \item DP with a bounded jump length.
  \item Minimum-cost path with moves limited to the last \(k\) states.
  \item Sequence partition problems reducible to sliding minima.
\end{itemize}

\section*{References}

\begin{itemize}
  \item \href{https://codeforces.com/catalog}{Codeforces Catalog}.
  \item \href{https://usaco.guide/plat/convex-hull-trick}{USACO Guide sections on DP optimization for related patterns}.
\end{itemize}
